\documentclass[11pt,a4paper]{article}
\usepackage[margin=1in]{geometry}
\usepackage{amsmath,amssymb}
\usepackage{booktabs}
\usepackage{array}
\usepackage{listings}
\usepackage{xcolor}
\definecolor{seclink}{HTML}{1F4E8C}
\usepackage[colorlinks=true,linkcolor=seclink,citecolor=seclink,urlcolor=blue]{hyperref}
\usepackage{parskip}
\usepackage{tikz}
\usetikzlibrary{positioning,arrows.meta,fit,backgrounds,calc}

\definecolor{ink}{HTML}{1A1A1A}
\definecolor{grey}{HTML}{5F5B54}
\definecolor{pale}{HTML}{F5F3EE}
\definecolor{tensor}{HTML}{1F6F4A}
\definecolor{seq}{HTML}{1F4E8C}
\definecolor{prod}{HTML}{9C3B1B}
\definecolor{hdrbg}{HTML}{D6E4F5}   % headers
\definecolor{bdybg}{HTML}{D5EDDC}   % bodies
\definecolor{linebg}{HTML}{F7D8E3}   % protocol / request line / status line
\definecolor{pathbg}{HTML}{F6E7C1}   % paths

\tikzset{
  srv/.style={draw=ink, fill=white, rounded corners=2pt, inner sep=3.4pt,
              font=\scriptsize\ttfamily, minimum height=5.6mm},
  plan/.style={srv, fill=pale, font=\scriptsize\ttfamily\bfseries},
  op/.style={circle, draw, inner sep=0.6pt, minimum size=6.4mm, font=\footnotesize\bfseries},
  ar/.style={-{Stealth[length=1.8mm]}, semithick, color=grey},
  lbl/.style={font=\tiny\itshape, color=grey},
  capt/.style={font=\scriptsize\bfseries},
  subt/.style={font=\tiny\itshape, color=grey},
}

\definecolor{kw}{RGB}{20,80,160}
\definecolor{cmt}{RGB}{110,110,110}
\definecolor{str}{RGB}{160,60,20}
\definecolor{codebg}{RGB}{246,246,244}

\lstdefinestyle{code}{
  basicstyle=\ttfamily\small,
  keywordstyle=\color{kw}\bfseries,
  commentstyle=\color{cmt}\itshape,
  stringstyle=\color{str},
  breaklines=true,
  frame=leftline,
  framesep=7pt,
  xleftmargin=12pt,
  rulecolor=\color{gray!45},
  backgroundcolor=\color{codebg},
  columns=fullflexible,
  showstringspaces=false,
  keepspaces=true,
  literate=
    {↓}{{$\downarrow$}}1 {↑}{{$\uparrow$}}1
    {→}{{$\rightarrow$}}1 {⊠}{{$\boxtimes$}}1
    {◁}{{$\lhd$}}1 {×}{{$\times$}}1 {—}{{\textemdash}}1 {’}{{'}}1,
}
\lstset{style=code}

% lstlinebgrd, patched for current listings (the 'Numbers none unknown' bug).
\makeatletter
\let\old@lstKV@SwitchCases\lstKV@SwitchCases
\def\lstKV@SwitchCases#1#2#3{}
\makeatother
\usepackage{lstlinebgrd}
\makeatletter
\let\lstKV@SwitchCases\old@lstKV@SwitchCases
\lst@Key{numbers}{none}{%
  \def\lst@PlaceNumber{\lst@linebgrd}%
  \lstKV@SwitchCases{#1}%
  {none:\\%
   left:\def\lst@PlaceNumber{\llap{\normalfont
     \lst@numberstyle{\thelstnumber}\kern\lst@numbersep}\lst@linebgrd}\\%
   right:\def\lst@PlaceNumber{\rlap{\normalfont
     \kern\linewidth \kern\lst@numbersep
     \lst@numberstyle{\thelstnumber}}\lst@linebgrd}%
  }{\PackageError{Listings}{Numbers #1 unknown}\@ehc}}
\makeatother

% Enforcement badges: compile time / run time / developer discipline.
\definecolor{enfc}{HTML}{1F4E8C}
\definecolor{enfr}{HTML}{1F6F4A}
\definecolor{enfd}{HTML}{9C3B1B}
\newcommand{\badge}[2]{\raisebox{0.05ex}{\fcolorbox{#1}{white}{\textcolor{#1}{\scriptsize\bfseries\,#2\,}}}}
\newcommand{\enfC}{\badge{enfc}{C}}
\newcommand{\enfR}{\badge{enfr}{R}}
\newcommand{\enfD}{\badge{enfd}{D}}

% Keep every code block whole: wrapping each listing in an unbreakable minipage
% makes LaTeX move it to the next page rather than split it across one.
\usepackage[most]{tcolorbox}
\definecolor{yellowbg}{HTML}{FCF6DC}
\tcbset{plainbg/.style={breakable, enhanced, sharp corners, boxrule=0pt, frame hidden,
  parbox=false,
  before upper={\setlength{\parskip}{0.55\baselineskip}\setlength{\parindent}{0pt}},
  left=4mm, right=4mm, top=3mm, bottom=3mm, before skip=3mm, after skip=3mm}}
\newtcolorbox{listbox}{plainbg, colback=black!5}
\newtcolorbox{orestisbox}{plainbg, colback=hdrbg}
\newtcolorbox{sectionbox}{plainbg, colback=yellowbg}
\newtcolorbox{pinkbox}{plainbg, colback=linebg}

\usepackage{etoolbox}
\BeforeBeginEnvironment{lstlisting}{\par\vspace{0.4\baselineskip}\noindent\begin{minipage}{\linewidth}}
\AfterEndEnvironment{lstlisting}{\end{minipage}\par\vspace{0.4\baselineskip}}

% Terminal trace: literate arrows are mis-placed under columns=fullflexible,
% and a literate glyph opening line 1 is flushed onto its own line -- hence
% keepspaces and the leading space on every line of the block below.
\lstdefinestyle{trace}{
  basicstyle=\ttfamily\small,
  keepspaces=true,
  frame=leftline, framesep=7pt, xleftmargin=12pt,
  rulecolor=\color{gray!45},
  backgroundcolor=\color{codebg},
  literate={↓}{{$\downarrow$}}1 {↑}{{$\uparrow$}}1 {…}{{\ldots}}1,
}

\lstdefinelanguage{TS}{
  morekeywords={type,function,switch,case,return,const,extends,never,interface,let,await,async,export,import,null,unknown,new,try,catch,throw,in,of},
  morecomment=[l]{//},
  morestring=[b]",
}
\lstdefinelanguage{SH}{
  morekeywords={curl,deno,run,export,for,do,done,if,then,fi,echo,sort,head,chmod,cat},
  morecomment=[l]{\#},
  morestring=[b]',
}

\title{\vspace{-2em}TypeScript Syntax}
\author{Robin Langer}
\date{1 September 2026}

\begin{document}
\maketitle

\begin{center}
\begin{minipage}{0.88\linewidth}\small\itshape
\textbf{Companion.} This is the TypeScript reference for the note \emph{Stalin: The Game}.
It is not meant to be read straight through. Refer back to it as necessary.
\end{minipage}
\end{center}
\vspace{1mm}

\tableofcontents

\newpage

\section{TypeScript syntax}
\label{sec:ts}

\begin{sectionbox}
It is an unfortunate accident of history that JavaScript became the ``default'' language for web
app development. TypeScript is an unfortunate attempt to ``bolt on'' a type system to JavaScript
after the event.

This is a ``crash course'' in TypeScript, and I am also learning as I write here. It is a bit of a
head twister, and it bears almost no resemblance to Haskell or Lean.

Don't try to memorize all this now. Refer back to it as necessary.
\end{sectionbox}

\begin{center}
\colorbox{hdrbg}{\begin{minipage}{0.94\linewidth}\vspace{1.5mm}
\textbf{The system is deliberately unsound.} This is the sentence to carry into the rest of the
note. TypeScript does not claim that a well-typed program cannot go wrong. It claims to catch a
useful class of mistakes.
\vspace{1.5mm}\end{minipage}}
\end{center}

\subsection{Basic types}
\label{sec:tssyntax}

\begin{listbox}
\begin{itemize}\setlength{\itemsep}{2pt}

\item The atomic types are \texttt{number}, \texttt{string}, \texttt{boolean}, \texttt{null} and
      \texttt{undefined}.

\item A string in quotes may be used as a type. The type \texttt{"failure"} has exactly one
      inhabitant, namely that one string.

\item The type \texttt{never} has no inhabitants at all.

\item The type \texttt{A | B} is a value that is an \texttt{A} or a \texttt{B}. It is an
      untagged union.

\item The type \texttt{\{ id: number; title: string; author: string \}} is a \emph{record}. A
      value of it is a dictionary in the Python sense, with those three keys. In our running
      example this is the type of a book, and the note gives it the name \texttt{Book}.

\item The type \texttt{Array<Book>} is a list whose elements are all books.

\item The signature \texttt{<T, R>(x: T) => R} is a generic function.

\item The keyword \texttt{export} in front of a definition means that other files may use it.
      Without it, the definition is private to its own file.

\end{itemize}
\end{listbox}

\subsection{The weird stuff}

\begin{listbox}
\begin{itemize}\setlength{\itemsep}{2pt}

\item The type \texttt{unknown} is a value about which nothing is known. Nothing may be done to
      it until you have established what it is.

\item The type \texttt{Record<string, unknown>} is a dictionary again, but this time with
      arbitrary keys and nothing at all known about the values.

\item The keyword \texttt{any} switches the type checker off for that value entirely.

\item The notation \texttt{u as T} says ``trust me, this is a \texttt{T}''. This is very
      upsetting for somebody in formal verification. It is called a \emph{type assertion}.

\end{itemize}
\end{listbox}

\newpage

\subsection{Type guards}
\label{sec:guards}

The notation \texttt{u as T} above asserts a type and checks nothing. There is a second notation
that checks, and it is the one the wire code in this note is built on.

\begin{lstlisting}[language=TS]
export const isRecord = (u: unknown): u is Record<string, unknown> =>
  typeof u === "object" && u !== null && !Array.isArray(u);
\end{lstlisting}

The word \texttt{is} is a keyword. It may appear only in the return type of a function, written
\texttt{u is T} where \texttt{u} is one of that function's own parameters. It says that the function
returns a boolean, and that a true answer means \texttt{u} is a \texttt{T}. A function declared that
way is called a \emph{type guard}.

The type guard above is used in one place only, immediately after a \texttt{JSON.parse}. So the
value it is asked about is never an arbitrary TypeScript value. It is one of the six things the
JSON grammar of the previous section allows, and the test only has to tell those six apart.

The body of \texttt{isRecord} is three tests joined by \texttt{\&\&}. The operator \texttt{typeof}
answers with a string naming the kind of a value.

The first test, \texttt{typeof u === "object"}, rules out numbers, strings and booleans. The
second, \texttt{u !== null}, rules out \texttt{null}. The third, \texttt{!Array.isArray(u)}, rules
out arrays, which \texttt{typeof} also calls objects. This leaves
\texttt{Record<string, unknown>} as the only remaining option.

\newpage

\subsection{Subtypes}

\begin{listbox}
\begin{itemize}\setlength{\itemsep}{2pt}

\item In TypeScript we can talk about one type being a \emph{subtype} of another. For example
      the type \texttt{"failure"} is a subtype of \texttt{string}. Every type is a subtype of
      \texttt{unknown}, and \texttt{never} is a subtype of every other type.

\item The type \texttt{Extract<Q, S>} keeps only those members of the union \texttt{Q} that are
      subtypes of \texttt{S}, and throws the rest away.

\item The notation \texttt{Book[K]} tells you the type of the value stored at field \texttt{K}
      of the record type \texttt{Book}. It is called an \emph{indexed access type}. So
      \texttt{Book["id"]} is \texttt{number}, and \texttt{Book["title"]} is \texttt{string}.

\item The constraint \texttt{K extends U} says that \texttt{K} must be a subtype of \texttt{U}.

\end{itemize}
\end{listbox}

\subsection{The powerful stuff}

\begin{listbox}
\begin{itemize}\setlength{\itemsep}{2pt}

\item The notation \texttt{X extends Y ? A : B} is a \emph{conditional type}. It computes a type
      by testing \texttt{X} against \texttt{Y}. This is the construct everything else in this
      note rests on, and it gets a subsection of its own below.

\item The notation \texttt{\{ [K in U]: F<K> \}} is a \emph{mapped type}, with one field for
      every member of the union \texttt{U}. This one also gets a subsection of its own below.

\end{itemize}
\end{listbox}

\subsection{On the level of values}

\begin{listbox}
\begin{itemize}\setlength{\itemsep}{2pt}

\item The keywords \texttt{interface} and \texttt{type} both name a type, the first restricted to
      records.

\item The declaration \texttt{const x: T = e} introduces a constant called \texttt{x}, gives it
      the type \texttt{T}, and sets it equal to the value \texttt{e}. The type may be left out.
      Writing \texttt{const n = 3} gives \texttt{n} the type \texttt{number}, because \texttt{3}
      is a number and the checker can see that.

\item The expression \texttt{(x) => e} is a lambda expression.

\item The expression \texttt{x ?? y} is equal to $x$, unless $x$ is \texttt{null} or
      \texttt{undefined}, in which case it is equal to $y$.

\item The expression \texttt{c ? a : b} is the ordinary conditional expression on the level of
      values.

\end{itemize}
\end{listbox}

\newpage

\section{Structural typing}
\label{sec:structural}

\begin{sectionbox}
A type describes what fields a value has. Any value with those fields has that type, whether or
not anybody said so. Nothing has to be declared and nothing has to be constructed.
\end{sectionbox}

For unions that has a cost, and an example makes it plain. Suppose a length is either in metres
or in feet.

\begin{lstlisting}[language=TS]
type Metres = { value: number };
type Feet   = { value: number };
type Length = Metres | Feet;
\end{lstlisting}

Those two are the same type. A value with one numeric field called \texttt{value} is a
\texttt{Metres}, and it is equally a \texttt{Feet}, and the checker will let you pass either one
where the other is expected. Given a \texttt{Length} there is no test that tells you which you
have, because there is nothing to test.

In Haskell you would write \texttt{data Length = Metres Double | Feet Double}. The constructor is
itself the answer. TypeScript has no constructors, so if you want to tell the two apart you must
put something there to look at.

\begin{lstlisting}[language=TS]
type Metres = { unit: "m";  value: number };
type Feet   = { unit: "ft"; value: number };
\end{lstlisting}

Now \texttt{x.unit === "m"} answers the question. But notice what \texttt{unit} is. It is an
ordinary field holding an ordinary string. The language did not provide it. We added it, and we
have to keep adding it, and everything that wants to tell metres from feet has to read it.

\subsection{A tagged union writes down a finite choice}

Once every case carries a tag, the checker can use it. Stay with lengths, and give the two cases
different data so that there is something to get wrong.

\begin{lstlisting}[language=TS]
type Length =
  | { unit: "m";  metres: number }
  | { unit: "ft"; feet: number };
\end{lstlisting}

\newpage

\section{Compiling down to JavaScript}
\label{sec:unusual}

\subsection{A cast is a promise, not a check}

The types are \emph{erased}. The program that executes has no representation of them and no check
derived from them. So the \texttt{as} operator is not a conversion. It is an assertion. The checker
accepts it without verifying, and the runtime cannot verify it at all.

In the example below, \texttt{toFixed} is a method that rounds a number to a given number of
decimal places. It exists on numbers and on nothing else, which is what makes it a useful test.

\begin{lstlisting}[language=TS]
const raw = '{"id":"1","title":"Diaspora","author":"Greg Egan"}';
const b = JSON.parse(raw) as Book;   // a promise, not a check
console.log(b.id.toFixed(0));        // compiles: id is declared number
\end{lstlisting}

\begin{lstlisting}[language=SH]
$ deno check t2.ts      -- passes. Nothing is wrong as far as the checker knows.
$ deno run   t2.ts
TypeError: b.id.toFixed is not a function
\end{lstlisting}

The checker was told \texttt{id} is a number and believed it. The id was a string all along. The
error surfaced at the first place that actually used it. That may be a long way from the cast, and
in a real program is likely to be somewhere else entirely.

\subsection{\texttt{unknown} is what you want at a boundary}

Two types describe a value you know nothing about, and they behave in opposite ways.

The type \texttt{any} switches the checker off. Every operation on an \texttt{any} is permitted.

\begin{lstlisting}[language=TS]
const raw: any = JSON.parse('{"id":"1","title":"Diaspora","author":"Greg Egan"}');
console.log(raw.id.toFixed(0));
\end{lstlisting}

\begin{lstlisting}[language=SH]
$ deno check anyex.ts     -- passes, no complaint
$ deno run   anyex.ts
TypeError: raw.id.toFixed is not a function
\end{lstlisting}

The id was the string \texttt{"1"} all along. Nothing objected, because nothing was looking.

The type \texttt{unknown} does the opposite. It permits nothing until you have established what you
have.

\begin{lstlisting}[language=TS]
const raw: unknown = JSON.parse('{"id":"1","title":"Diaspora","author":"Greg Egan"}');
console.log(raw.id);
\end{lstlisting}

\begin{lstlisting}[language=SH]
$ deno check unknownex.ts
TS18046 [ERROR]: 'raw' is of type 'unknown'.
\end{lstlisting}

That refusal is the useful thing. The type \texttt{any} would have compiled and then failed later, in the
manner of the previous subsection. The type \texttt{unknown} makes the boundary visible. It refuses to let
you cross it accidentally.

\subsection{Validate, do not cast}

So how does a value we know nothing about honestly become one whose type we know? By code that looks at
the fields. The signature to aim for is \texttt{unknown} in, and the type \emph{or nothing} out.

\begin{lstlisting}[language=TS]
function asBook(u: unknown): Book | null {
  if (typeof u !== "object" || u === null) return null;
  const o = u as Record<string, unknown>;
  if (typeof o.id     !== "number") return null;
  if (typeof o.title  !== "string") return null;
  if (typeof o.author !== "string") return null;
  return { id: o.id, title: o.title, author: o.author };  // rebuilt, not asserted
}
\end{lstlisting}

\newpage

Try it on three documents.

\begin{lstlisting}[language=TS]
asBook({ id: 13,   title: "Diaspora", author: "Greg Egan" })
// { id: 13, title: "Diaspora", author: "Greg Egan" }

asBook({ id: "13", title: "Diaspora", author: "Greg Egan" })
// null -- the id is a string

asBook({ id: 13,                     author: "Greg Egan" })
// null -- there is no title
\end{lstlisting}

\end{document}
